% !TeX root = ../main.tex

\newlecture[Probability Measure][Sep 10, 2026]

\deflabel{2.1}{Sample Space}
\begin{quote}
    A sample space $\Omega$ is any non-empty set.
\end{quote}

\begin{example}
    The sample space contains ''all the states of the following''
    \\Coin Flipping: $\Omega: \{H, T\}$
    \\Weather: $\{\text{hot, cold, mild}\} \times \{\text{wet, dry}\}$
\end{example}

\deflabel{2.2}{Probability Measure (STA347 ver.)}
\begin{quote}
    For a sample space $\Omega$, we say $P$ is a probability measure if the following holds:
    \begin{enumerate}[label=\roman*)]
        \item $\forall E \subseteq \Omega, 0 \leq P(E) \leq 1$, that is $P:\mathcal{P}(\Omega) \to [0,1]$
        \item $P(\Omega) = 1$
        \item $\forall E_1, E_2, \dots \subseteq \Omega$ such that $E_i \cap E_j = \varnothing \ \forall\ i \neq j$
        \[P\left(\bigcup_{i=1}^{\infty} E_i\right) = \sum_{i=1}^{\infty}P(E_i)\]
    \end{enumerate}
\end{quote}

\aside{$\mathcal{P}(\Omega)$ is the power set of $\Omega$.}

\textbf{Remark:} It may be impossible for such a $P$ to exist. To define exactly which subsets we want $P$ to satisfy these axioms, that would require more advanced mathematical theories. For this class, we just assume such a $P$ is defined for all subsets of $\Omega$.

\begin{example}
    Probabilities of coin flipping
    \[\Omega = \{H, T\}\]
    \[\mathcal{P}(\Omega) = \{\phi, \{H\}, \{T\}, \{H, T\}\}\]
    \[P: \mathcal{P}(\Omega) \mapsto [0, 1]\]
    \begin{align*}
        p: \phi &\mapsto 0 \\
        \{H\} &\mapsto \frac{1}{2} \\
        \{T\} &\mapsto \frac{1}{2} \\
        \{H, T\} &\mapsto 1
    \end{align*}
\end{example}

\begin{example}
    The uniform/Lebesgue measure.
    Let $\Omega =[L,U]$ for some $L < U \in \R$. Define
    \[P([a,b]) = \frac{b-a}{U-L} \quad \text{when } L \leq a \leq b \leq U\]
    We assume such a measure exists (and is well defined) for this class.
\end{example}

Proposition 2.3:
The following properties are satisfied by any probability measure $P$
\begin{enumerate}
    \item $\forall E \subseteq \Omega, P(E^C) = 1 - P(E)$
    \item $P(\phi) = 0$
    \item $\forall E, F \subseteq \Omega$, such that $E \subseteq F, P(E) \leq P(F)$ \aside{(``big event, big probability'')}
    \item $\forall E, F \subseteq \Omega, P(E\bigcup F) = P(E)+P(F) - P(E\bigcap F)$
    \item $\forall E, F \subseteq \Omega, P(E\bigcap F) = P(E)+P(F)-P(E\bigcup F)$
    \item $\forall E, F \subset \Omega, P(F \setminus E) = P(F) - P(E\bigcap F)$
\end{enumerate}

Proof via Def 2.2 (probability axioms).

Lemma 2.4
\begin{quote}
    For a sequence $E_1, \dots, E_n \subseteq \Omega$
    \[P(\bigcup_{i=1}^n E_i) \leq \sum_{i=1}^{n}P(E_i)\]
    Proof: We define $F_i = E_i \setminus \bigcup_{j=1}^{i-1}E_j, F_1 = E_1$
    \[\bigcup_{i=1}^{n} E_i = \bigcup_{i=1}^{n}F_i \text{ and } F_i \bigcap F_j = \phi \text{ if } i \neq j\]
    from the probability axiom.
    \begin{align*}
        \text{Since } \bigcup_{i=1}^{n} E_i &= \bigcup_{i=1}^{n}F_i\\
        P(\bigcup_{i=1}^{n} E_i) &= P(\bigcup_{i=1}^{n}F_i)\\
        &= \sum_{i=1}^{n}P(F_i)
    \end{align*}
    \[\text{We also know } F_i \subseteq E_i \text{ as } F_i = E_i \setminus \bigcup_{j=1}^{i-1}E_j\]
    \[P(F_i) \leq P(E_i)\]
    \[P(\bigcup_{i=1}^{n}E_i) = \sum_{i=1}^{n}P(F_i) \leq \sum_{i=1}^{n}P(E_i)\]
\end{quote}

Lemma 2.5
\begin{quote}
    Let $P$ be the Lebesgue measure on $[0,1]$. Then $P(E) = 0$ for any \underline{countable} set $E$.
\end{quote}
